% Generated draft: requires mathematical and source review.
\documentclass[11pt]{article}
\usepackage[margin=25mm]{geometry}
\usepackage{fontspec}
\setmainfont{Latin Modern Roman}
\usepackage{amsmath,amssymb,mathtools,unicode-math}
\setmathfont{Latin Modern Math}
\usepackage{enumitem,needspace,xurl,hyperref}
\setlength{\parindent}{0pt}
\setlength{\parskip}{5pt}
\setlength{\emergencystretch}{3em}
\newcommand{\N}{\mathbb{N}}
\newcommand{\Z}{\mathbb{Z}}
\newcommand{\Q}{\mathbb{Q}}
\newcommand{\R}{\mathbb{R}}
\newcommand{\C}{\mathbb{C}}
\newcommand{\F}{\mathbb{F}}
\newcommand{\A}{\mathbb{A}}
\newcommand{\PP}{\mathbb{P}}
\newcommand{\E}{\mathbb{E}}
\newcommand{\eps}{\varepsilon}
\newcommand{\dd}{\,\mathrm{d}}
\begin{document}
\section*{UnsolvedMath ID: 30005684}
\textbf{Source problem code:} catalogue entry
\subsection{Definitions and mathematical statement}
Let \(\nu\) be a compactly supported probability measure on \(\R\), let \(S=\operatorname{supp}\nu\) be its topological support, and assume \(\#S\geq 2\). Set
\[
\Omega=S^{\Z},\qquad W_\omega(n)=\omega_n\quad(\omega\in\Omega,\ n\in\Z).
\]
Thus the noise is indexed by the entire integer line. Its multiplication operator on \(\ell^2(\Z)\) is
\[
(M_\omega u)(n)=\omega_nu(n).
\]
Compactness of \(S\subset\R\) makes this a bounded self-adjoint operator. If an unperturbed self-adjoint operator \(H_x\) on \(\ell^2(\Z)\) is specified, its additive perturbation can therefore be written conditionally as
\[
H_{x,\omega}=H_x+M_\omega,\qquad
\mathcal D(H_{x,\omega})=\mathcal D(H_x).
\]
This conditional notation supplies an interpretation of adding the displayed noise; it does not recover the missing definition of \(H_x\).

The excerpt describes an ergodic-potential setting and states that there is a compact set \(\Sigma_0\) such that
\[
\sigma(H_x)=\Sigma_0\qquad\text{for }\mu\text{-almost every }x\in X.
\]
However, it does not supply the definitions or full hypotheses for \(X\), \(\mu\), the base transformation \(T\), the sampling function, or \(H_x\). This almost-sure spectral assertion is consequently retained as supplied context, without asserting that it follows from hypotheses fully recorded here. The law of \(\omega\) on \(\Omega\), its independence properties, and its relation to the base variable \(x\) are also not explicitly specified.

The source's question is what happens when random noise is added to an aperiodic potential, or more generally to a given ergodic potential. Its examples include periodic potentials with \(X\) finite and \(T\) a cyclic permutation; almost Mathieu potentials with \(X=S^1\), \(T\) an irrational rotation, and \(f(x)=2\lambda\cos(2\pi x)\); Fibonacci substitution potentials; and Sturmian potentials with a Cantor base. The excerpt gives no parameter range for \(\lambda\). It also says that one-dimensional aperiodic structures can be included through sequence hulls, but does not specify the required invariant measure or ergodicity assumptions.

\textbf{Unresolved formulation.} No precise conjecture or requested conclusion can be recovered from this record. The excerpt ends at \emph{It turns out that}, before the subsequent conclusions. The supplied verification field labels the statement unrecoverable and reports that the introductory question is followed by localization and spectral-structure theorems, with a remaining Fibonacci case recorded separately. That separate formulation is not supplied here. Accordingly, this draft records the noise construction and the introductory question, without turning them into an asserted open conjecture.
\subsection{Short English statement}
The record asks about adding compactly supported random noise to an aperiodic or ergodic Schr\"odinger potential. It specifies the noise coordinates, but omits the full operator setup, probability law, and intended conclusion. The supplied verification labels this excerpt unrecoverable; a precise open problem cannot be reconstructed from it.
\subsection{Sources}
\begin{itemize}
\item Supplied local UnsolvedMath registry record, permanent ID \(30005684\), title \emph{Random Perturbations of Aperiodic Schr\"odinger Potentials}, exact source code \texttt{catalogue entry}, in the Oberwolfach Reports Open Problems collection. The record supplies an \emph{open} status label, but also \texttt{statement\_status: unrecoverable}; neither constitutes independent mathematical verification.
\item Supplied report citation: \emph{Aspects of Aperiodic Order (2024)}, Oberwolfach Reports, pp.\ 2107--2178, DOI \texttt{10.4171/owr/2023/37}. \url{https://doi.org/10.4171/owr/2023/37}. The record gives proposed year 2023. Its excerpt attributes the presentation \emph{On spectrum of aperiodic Schr\"odinger operators with random noise} to Anton Gorodetski, jointly with A.\ Avila, D.\ Damanik, and V.\ Kleptsyn; no separate proposer is identified.
\item Supplied literature source labeled Avila--Damanik--Gorodetski: \url{https://arxiv.org/abs/2211.02173}.
\item Supplied literature source labeled Damanik--Gorodetski--Kleptsyn: \url{https://arxiv.org/abs/2504.08153}.
\item Literature assessment dated exactly \texttt{2026-08-22}: the supplied curation labels the topic open and states that randomly perturbed ergodic potentials with connected base are understood and have interval spectrum, while the general aperiodic-noise question and disconnected Fibonacci hull remain unresolved. This is supplied dated curation derived from the Oberwolfach record, not an independently verified literature assessment; the hypotheses and scope of those claims require checking in the cited sources.
\item The supplied statement-verification note says that the official report answers the introductory question through localization and spectral-structure theorems and preserves the remaining Fibonacci case in a following record. This description is supplied verification metadata; the report and following record have not been independently inspected here.
\end{itemize}
\end{document}
